\section{L'équité, la diversité et l'inclusion}\label{lequite-la-diversite-et-linclusion}

\stanceinfo{\textbf{Congrès de la FCEG 2019} [2019-01-04]}{2024-01-02}

\officialstance{La Fédération canadienne étudiante de génie (FCEG) estime que l'équité, la diversité et l'inclusion (EDI) sont essentielles à la durabilité de la profession d'ingénierie et à sa contribution à l'intérêt public. L'inclusion et la participation accrue des groupes sous-représentés apportent des perspectives et des expériences différentes qui aident les ingénieurs à résoudre les problèmes au sein de la profession d'ingénieur et dans les programmes d'études de premier cycle.}

\stanceheading{La position des étudiants}
\begin{itemize}
 \item Le renforcement de la diversité et de la représentation des groupes minoritaires au sein de l'ingénierie permet de relever les défis et de créer des solutions aux problèmes d'ingénierie qui profitent à tous.
 \item Les étudiants en ingénierie doivent avoir accès à l'égalité des chances dans le domaine de l'ingénierie, indépendamment de leur âge, de leur origine ethnique, de leur race, de leur sexe, de leur orientation sexuelle, de leur situation économique, de leur capacité, de leur religion ou d'autres origines diverses. Ces possibilités comprennent l'accès aux ressources et aux possibilités qui favorisent leur développement en tant qu'ingénieur (stages, conférences sur l'ingénierie, etc.).
 \item L'augmentation de la diversité des étudiants dans les programmes d'ingénierie permet d'améliorer la formation des ingénieurs en exposant les étudiants à davantage de perspectives, d'expériences et de manières de penser.
 \item La diversité sera favorisée à un taux plus élevé et plus efficace dans un environnement inclusif en contribuant à la rétention des groupes sous-représentés. Afin de créer un environnement inclusif, il convient de mettre en place une éducation et des initiatives autour des thèmes de la diversité afin d'éduquer les individus à améliorer l'accessibilité pour les personnes diverses et à réduire les préjugés inconscients.
 \item L'intégration des populations autochtones dans l'ingénierie est un sujet important, mais peu discuté, étant donné que le Canada se trouve sur des terres autochtones traditionnelles et non cédées et que son existence est le résultat d'une colonisation dont les effets sont encore présents aujourd'hui.
 \item La FCEG approuve également les Énoncés de principe nationaux d'Ingénieurs Canada et l'initiative 30 en 30 d'Ingénieurs Canada, qui vise à accroître la proportion des personnes s'identifiant en tant qu'ingénieurs nouvellement agréés à 30 pour cent d'ici 2030.
 \item Les groupes sous-représentés sont souvent désavantagés en matière d'accès à l'ingénierie, d'accessibilité et de capacité à prospérer en raison d'obstacles systémiques qui les affectent de manière disproportionnée.
 \item Créer un environnement inclusif au sein de l'ingénierie implique l'établissement d'un cadre, d'une culture et de pratiques sécurisés qui accueillent et respectent les contributions des individus issus de milieux, d'expériences et de perspectives divers.
\end{itemize}

\stanceheading{Les recommandations aux partenaires, aux parties intéressées et aux autres entités}
\begin{itemize}
 \item La FCEG demande au Réseau 30 en 30 d'Ingénieurs Canada d'étudier les raisons pour lesquelles les taux de rétention des femmes sont faibles dans les études de premier cycle et de les présenter officiellement à d'autres parties prenantes, y compris au conseil d'administration d'Ingénieurs Canada.
 \item La FCEG demande au Réseau 30 en 30 d'Ingénieurs Canada de commencer à élaborer des plans secondaires si le programme 30 en 30 n'est pas réalisé à l'échelle du Canada, et d'intégrer les voix des étudiants dans la nouvelle proposition.
 \item La FCEG demande aux DDIC de continuer à promouvoir le recrutement d'identités diverses parmi les nouveaux étudiants en génie au moyen d'initiatives qui favorisent l'équité au sein de la communauté des ingénieurs. (i.e. bourses d'études pour les communautés sous-représentées)
 \item La FCEG demande au BCAPG d'intégrer l'équité, la diversité et l'inclusion (ÉDI) dans les attributs des diplômés, de façon similaire à l'impact du génie sur la société et l'environnement, de décrire l'importance de la formation en ÉDI et de toujours s'assurer que les considérations liées à l'ÉDI sont prises en compte dans la profession de génie.
 \item La FCEG demande aux DDIC d'intégrer une vice-doyenne à l'équité, à la diversité et à l'inclusion ou un poste de rang similaire au sein de leur établissement afin d'assurer une accessibilité équitable, diversifiée et inclusive au génie et d'offrir un soutien aux groupes sous-représentés.
 \item La FCEG demande aux DDIC d'offrir une formation sur l'équité, la diversité et l'inclusion à leurs instructeurs et à leurs étudiants, en particulier aux étudiants qui représentent l'ensemble de l'établissement.
 \item La FCEG demande aux administrateurs des établissements, avec le soutien des doyens des facultés d'ingénierie du Canada, de mettre en œuvre et d'appliquer des protocoles pour traiter les cas de discrimination et de harcèlement au sein de la communauté des ingénieurs.
 \item La FCEG incite le BCAPG et les DDIC à promouvoir les principes de gestion dans les programmes d'études, afin que les étudiants soient mieux équipés pour prendre en compte et défendre les valeurs de l'ÉDI sur leur lieu de travail à la fin de leurs études.
 \item Enfin, la FCEG incite nos membres à examiner les traditions d'ingénierie qu'ils défendent et à considérer avec un regard critique si elles respectent et contribuent à créer un environnement inclusif pour tous. En outre, la FCEG incite ses membres à rechercher des occasions de promouvoir la connaissance d'ÉDI parmi les étudiants et à créer des initiatives qui soutiennent l'équité, la diversité et l'inclusion au sein de leurs institutions.
\end{itemize}

\stanceheading{Ce que la FCEG envisage de faire}

La FCEG prévoit de renforcer son soutien et sa promotion de l'inclusion au sein de la communauté des étudiants en ingénierie en fournissant des ressources et en partageant les meilleures pratiques avec les écoles membres en ce qui concerne l'organisation d'événements inclusifs et la mise en place d'orientations/de formations sur l'inclusion au sein de leurs communautés respectives.

\begin{itemize}
 \item Poursuivre ses efforts en matière d'inclusion en proposant des orientations sur l'inclusion lors de tous les événements, ainsi qu'en approfondissant l'évaluation des risques en matière d'inclusion lors des événements;
 \item Explorer de nouveaux partenariats avec des organisations axées sur la diversité et l'inclusion (Catalyst Canada, Actua, Women in Engineering, Association canadienne pour la santé mentale, National Society of Black Engineers, Laboratoire d'innovation en ingénierie, etc.);
 \item Donner la priorité à l'intégration active du bilinguisme dans toutes les activités de la FCEG, ainsi qu'au cœur de l'organisation;
 \item Élaborer des lignes directrices pour la planification d'événements inclusifs spécifiques aux activités de la FCEG afin d'aider à identifier, supprimer et prévenir les obstacles susceptibles d'empêcher quiconque de participer pleinement aux activités de la FCEG, et veiller à ce que ces considérations soient prises en compte dès les premières étapes de la planification de chaque activité;
 \item Respecter tous les engagements énoncés dans la Charte des champions afin d'appuyer officiellement l'initiative 30 en 30 d'Ingénieurs Canada, notamment en désignant chaque année, lors de la réunion de printemps, un champion qui participe au groupe des champions et qui s'engage activement dans le réseau 30 en 30;
 \item Mettre en pratique le formulaire de signalement des cas d'exclusion et des actes préjudiciables, qui permet aux membres de faire part de leurs commentaires;
 \item Donner la priorité aux événements et activités sans barrières lors des conférences de la FCEG, en proposant des alternatives lorsque cela n'est pas possible;
 \item Créer une boîte à outils pour les membres de la FCEG sur l'équité, la diversité et l'inclusion à l'intention des membres, qui comprendra notamment les éléments suivants :
 \begin{itemize}
  \item La formation sur l'équité, la diversité et l'inclusion
  \item Le programme de certification ÉDI
  \item Les pratiques exemplaires pour les organisations externes
  \item Les statistiques démographiques canadiennes sur le génie
  \item L'accès aux partenaires de la FCEG pour la défense des intérêts locaux
  \item Les mécanismes de retour d'information
 \end{itemize}
\end{itemize}

\newpage
